% 总论第 4 章 —— 验证方法学

\chapter{验证方法学}
\label{ch:verification-methodology}
\index{1concept@概念!verification@验证}

本章记述"凭什么相信这些数字"这件事本身。它不描述任何一个模块，
而是描述九个模块章里那些验证表\emph{共同遵循的判据体系}。
把它单列一章的理由很直接：本文档全篇出现的几百个数字，
其可信度完全取决于本章所定义的判据是否恰当。

\section{四层验证框架}
\label{sec:verify-framework}

方案把验证分为四层，本文档沿用其分类：

\begin{center}
\small
\begin{tabular}{@{}p{0.24\textwidth} p{0.34\textwidth} p{0.34\textwidth}@{}}
\toprule
\textbf{层} & \textbf{内容} & \textbf{本文档中的体现}\\
\midrule
代码验证 & 收敛阶、闭式解、制造解 & 各章 §x.5 的主体\\
解验证 & 网格收敛、容差与步长扫描 & 第 \ref{ch:sp1d}、\ref{ch:transport_dd} 章\\
确认（Validation） & 与实测/独立实现对比 & \textbf{仅部分完成}，见下\\
不确定度量化 & 敏感性、蒙特卡洛 & 未实现\\
\bottomrule
\end{tabular}
\end{center}

第三层是本文档最弱的一环，必须明说：\textbf{基准阶梯 L0--L7 目前只走到 L1}。
L2（Hall 对标面密度）及其后的每一级都尚未进行，因此
"绝对数值不可信"这句话在本文档中反复出现，其根据就在这里。

\begin{veriftable}{基准阶梯的当前进度}{tab:verify-ladder}
\vrow{L0\ \ 平衡态、一维 Poisson、极化解析式}{$<10^{-6}$}{\textbf{达成}（第 \ref{ch:sp1d}、\ref{ch:polarization} 章）}
\vrow{L1\ \ 一维 Schrödinger--Poisson 子带}{子带能量差 $<5$\U{meV}}{部分达成（无独立实现对照）}
\vrow{L2\ \ 标准异质结 $\ns$、C--V}{$\ns$ 差 $<5\%$}{\textbf{未做}（需霍尔测量）}
\vrow{L3\ \ 与商用 TCAD 交叉对比 IV/C--V}{$<5\%$}{未做}
\vrow{L4\ \ 直流输出/转移特性}{$I_{d\max}/g_m<10\%$，$V_{\text{th}}<0.15$\U{V}}{未做}
\vrow{L5\ \ S 参数 $\to f_T/f_{\max}$}{$<10$--$15\%$}{未实现（无 S 参数链）}
\vrow{L6\ \ pulsed IV、gate/drain-lag}{坍塌幅度 $<20\%$}{未实现（无陷阱模型）}
\vrow{L7\ \ load-pull}{$P_{\text{out}}$/PAE $<10$--$15\%$}{未实现}
\end{veriftable}

\section{五类判据}
\label{sec:verify-criteria}

本文档 §x.5 的验证表使用且只使用五类判据。每一类都有它的适用场景
与它\emph{不能}证明的东西 —— 知道后者比知道前者更重要。

\subsection{第一类：闭式解}
\label{sec:verify-closed-form}
\index{3method@方法!closed-form benchmark@闭式解基准}

\textbf{能证明}：离散格式在这个特定问题上的精度。
\textbf{不能证明}：格式在别的问题上也是对的。

闭式解是最强的单项证据，因为它是\emph{精确}的，比较不需要容差解释。
本文档中的典型值：

\begin{veriftable}{闭式解基准的代表性结果}{tab:verify-closed}
\vrow{双介质平行板（含界面面电荷）}{精确}{\textbf{$<10^{-10}$\U{V}}（\crefchap{ch:sp1d}）}
\vrow{无限深方阱 $E_n=n^2\pi^2\hb^2/(2\meff L^2)$}{二阶收敛}{二阶收敛（\crefchap{ch:sp1d}）}
\vrow{均匀电阻 $J=qN_D\mumu V/L$}{$<0.1\%$}{\textbf{$<10^{-6}$}（\crefchap{ch:transport_dd}）}
\vrow{块三对角解对稠密解}{$<10^{-10}$}{通过（四组尺寸，\crefchap{ch:blocktridiag}）}
\vrow{Bernoulli 函数 $B(x)=x/(\ee^x-1)$}{$10^{-12}$ 相对}{通过（\crefchap{ch:transport_dd}）}
\end{veriftable}

注意最后两行的性质差别：第一至三行检验的是\emph{离散格式}，
第四行检验的是\emph{求解器}，第五行检验的是\emph{一个特殊函数}。
把它们分开陈述，是因为它们失效时的含义完全不同。

\subsection{第二类：收敛阶}
\label{sec:verify-order}
\index{3method@方法!convergence order@收敛阶}

\textbf{能证明}：格式的\emph{阶}是对的 —— 这比"某个网格上误差小"
强得多，因为它意味着加密时误差会按预期速率下降。
\textbf{不能证明}：误差常数小。一个一阶格式可能在粗网格上比
二阶格式更准。

网格收敛阶由 Richardson 三点公式估计：

\begin{equation}
p \;=\;
\frac{\ln\!\bigl(\abs{u_1-u_2}\,/\,\abs{u_2-u_3}\bigr)}
     {\ln\!\bigl(h_1/h_2\bigr)},
\qquad h_1>h_2>h_3 .
\label{eq:verify-richardson}
\end{equation}

迭代法的\emph{局部收敛阶}由相邻三个残差估计：

\begin{equation}
p \;=\;
\frac{\ln\!\bigl(\norm{r_{k+1}}/\norm{r_k}\bigr)}
     {\ln\!\bigl(\norm{r_k}/\norm{r_{k-1}}\bigr)} .
\label{eq:verify-local-order}
\end{equation}

\crefeq{eq:verify-local-order} 的使用有一个必须说明的细节：
\textbf{只有落在渐近窗口内的三元组才有意义}。实现上取
"残差严格下降且比值本身在缩小"的三元组，并排除最小成员已落到
舍入地板（$\sim10^{-15}$）的样本 —— 地板以下的比值是噪声。
本文档报告的两个关键值都取自这样的窗口：耦合 SP 的
$p\approx1.99$（\crefchap{ch:sp1d}）与制造解的
$2.008/2.002/2.001$（\crefchap{ch:transport_dd}）。

\subsection{第三类：制造解}
\label{sec:verify-mms}
\index{3method@方法!method of manufactured solutions@制造解方法}

\textbf{能证明}：算子在\emph{任意}已知解上的精度，而不只是在恰好有
简单闭式解的那几个问题上。
\textbf{不能证明}：边界条件的处理是否正确（制造解通常搭配周期性或
简单边界）。

做法是：取一个解析函数 $\psi^{\star}$，把它代入微分算子算出源项，
再把该源项作为已知右端求解，比较数值解与 $\psi^{\star}$。
本文档中的范例在 \crefchap{ch:transport_dd}：取
$\psi^{\star}=\sin(\pi z/L)$（与横向无关，从而\emph{隔离}竖直算子），
加密四档得到的观测阶为 $2.008/2.002/2.001$。
标准表述见 \cite{roache2002,salari2000}。

制造解的一个常被忽略的好处是：它使\textbf{"界面以外的格式是二阶的"}
成为一个可断言的命题，即便界面本身只有一阶 —— 只要制造解在界面两侧
都是光滑的。\crefchap{ch:transport_dd} 正是这样把两件事分开陈述的。

\subsection{第四类：不变量与趋势}
\label{sec:verify-invariants}
\index{1concept@概念!invariant assertion@不变量断言}

\textbf{能证明}：某个\emph{必须成立的关系}在实现中确实成立。
\textbf{不能证明}：任何数值精度。

不变量断言的价值在于\emph{耐久}：它断言的是"若把 Al 组分调高，
极化必然增强"这类关系，而不是某个具体数字。
因此它在参数演化后依然有效 —— 这对一个参数尚未标定的项目
尤其重要。本文档中此类断言的例子：

\begin{itemize}[leftmargin=1.7em,itemsep=0.28em,topsep=0.35em]
  \item 平衡态通量恒为零（\crefchap{ch:transport_dd}，实测 $2.7\times10^{-15}$\U{A/m}）；
  \item 势偏导的反对称性 $\partial J/\partial\psi_L=-\partial J/\partial\psi_R$
        （规范不变性的离散形式）；
  \item 密度偏导与密度无关（跨 $17$ 个数量级验证）；
  \item 面密度低于极化上界 $\sigeff/q$；
  \item 屏蔽因子严格按比例缩放被测到 $10^{-6}$ 位。
\end{itemize}

\subsection{第五类：负向基准}
\label{sec:verify-negative}
\index{1concept@概念!negative benchmark@负向基准}

\textbf{能证明}：某个机制是\emph{必要}的。
\textbf{不能证明}：它是最优的。

负向基准记录的是\emph{另一个配置下的失败}。本文档中有两处：

\begin{veriftable}{负向基准（记录失败，以证明机制的必要性）}{tab:verify-negative}
\vrow{关闭线搜索后从远起点出发}{不收敛}{
  通过（$x^3-1$ 自 $x=40$、上限 $4$ 步；\crefchap{ch:newton}）}
\vrow{遗漏态密度质量响应时的误差地板}{不随差分步长下降}{
  $1.3\times10^{-5}$，是"缺项"的特征（\crefchap{ch:sp1d}）}
\vnote{第二行尤其值得注意：它记录的是一个\emph{已修正的}错误实现的症状。
      保留它，是因为"如何识别一个缺项"这件事无法用正面例子说清。}
\end{veriftable}

\section{导数校验：一个必须写进方法学的教训}
\label{sec:verify-jacobian}
\index{3method@方法!Jacobian verification@Jacobian 校验}

Newton 类方法的正确性几乎完全取决于 Jacobian。本文档记录了一条
\emph{代价高昂}的教训，它的普遍性超出了本项目。

\subsection{校验公式}

解析 Jacobian 与中心差分比较：

\begin{equation}
\mat{J}^{\text{FD}}_{:,j} \;=\;
\frac{\mat{R}(\bm{x}+h\bm{e}_j)-\mat{R}(\bm{x}-h\bm{e}_j)}{2h},
\qquad h = 10^{-6}\max\!\bigl(1,\norm{\bm{x}}_\infty\bigr).
\label{eq:verify-fd}
\end{equation}

分歧在于\textbf{如何把 $\mat{J}^{\text{FD}}_{:,j}-\mat{J}_{:,j}$ 归一化}。
两种做法：

\begin{equation}
e_j^{\text{col}} \;=\;
\frac{\bigl\lVert \mat{J}^{\text{FD}}_{:,j}-\mat{J}_{:,j}\bigr\rVert_\infty}
{\max\bigl(\bigl\lVert \mat{J}^{\text{FD}}_{:,j}\bigr\rVert_\infty,\
       \bigl\lVert \mat{J}_{:,j}\bigr\rVert_\infty\bigr)},
\qquad
e_i^{\text{row}} \;=\;
\frac{\bigl\lVert \mat{J}^{\text{FD}}_{i,:}-\mat{J}_{i,:}\bigr\rVert_\infty}
{\max\bigl(\bigl\lVert \mat{J}^{\text{FD}}_{i,:}\bigr\rVert_\infty,\
       \bigl\lVert \mat{J}_{i,:}\bigr\rVert_\infty\bigr)} .
\label{eq:verify-fd-norms}
\end{equation}

\subsection{为什么逐列不够}

\crefchap{ch:transport_dd} 的未知量从 $(\psi,\varphi_n)$ 改为 $(\psi,n)$ 时，
Jacobian 里残留了一行旧代码，把第二个未知量当作准费米势代进了态密度公式。
后果是\emph{内部密度与残差相差 $17$ 个数量级}，
连续性行因此完全失去对 $\psi$ 的耦合，整个耦合矩阵在 Newton 最需要的
方向上近奇异。

而当时的\emph{逐列}归一化校验报告相对误差只有 $10^{-10}$ ——
因为坏掉的那一项在整列中绝对值很小，被子列中量级正常的其他项淹没了。
改成\emph{逐行}相对判据后，它立刻暴露。

\boundary{由此得到两条纪律，它们适用于任何耦合系统的实现：

\textbf{（i）导数校验用逐行相对判据，不用逐列。}
行对应"一个方程"，列对应"一个未知量"；一个方程整体错位，
在行视角下是 $100\%$ 的错误，在列视角下可能微不足道。

\textbf{（ii）校验要在单元级与沿路径做两遍。}
"通量本身对中心差分"（单元级）与"整个 Jacobian 沿 Newton 路径
对中心差分"（系统级）是两道独立的闸门。\crefchap{ch:transport_dd}
两道都做了：前者逐项验四个偏导，后者沿路径验到 $\sim10^{-10}$。

\textbf{（iii）一致性不等于正确性。}
\crefchap{ch:transport_dd} 另记录了一例：Bernoulli 函数的反射恒等式
写反之后，通量仍然\emph{完全自洽} —— 系统级 Jacobian 对中心差分
照样通过 $4\times10^{-10}$ —— 却在每个带边沿电流方向上升的面上
系统性错误。修好后平衡态残差从 $35$ 降到 $4.8\times10^{-6}$（七个数量级）。
\textbf{自洽性只能排除"不自洽"，不能排除"一致地错"。}
}

\subsection{本文档中两个校验器的差异}

一个值得记录的细节：\crefchap{ch:newton} 的通用校验器用的是
\emph{逐列}归一化（\crefeq{eq:verify-fd-norms} 左式），
而 \crefchap{ch:transport_dd} 用自己的逐行判据。
这不是不一致，而是适用范围不同：前者的服务对象是良缩放的稠密系统，
后者是行均衡过的耦合系统。本文档在两处都注明了这一差异，
把它作为\emph{适用边界}而不是藏起来。

\section{可复现性作为一类判据}
\label{sec:verify-reproducibility}
\index{1concept@概念!reproducibility@可复现性}

有三类断言检验的不是精度，而是"同样的输入是否给出同样的东西"。
它们看似平凡，却是上述一切验证的前提。

\begin{veriftable}{三类可复现性断言}{tab:verify-repro}
\vrow{同输入同指纹}{两次运行的指纹相等}{
  通过（\crefchap{ch:sp1d} 的 $2\times$ 运行、\crefchap{ch:dft} 的同 Deck 两次运行）}
\vrow{换实现则指纹不同}{指纹不等}{
  通过（解析量子 Deck 与自洽量子 Deck、\crefchap{ch:dft} 的材料槽替换）}
\vrow{初值/延拓无关性}{解与初值路径无关}{
  \crefchap{ch:sp1d}：热启动 $0.35\to0.95$ 上界，相对差 $<0.5\%$；
  \crefchap{ch:transport_dd}：两条延拓阶梯，相对差 $<5\times10^{-3}$}
\vnote{第三行是本文档中最被低估的一类断言。它证明了"初值与延拓只影响
      收敛路径，不影响收敛解"—— 而这条性质一旦不成立，
      Newton 类方法的全部结果都会变得依赖调参历史，不再是方程的解。}
\end{veriftable}

\section{文档侧的表述纪律}
\label{sec:verify-doc-discipline}

最后，把上述判据体系翻译成一条写作纪律。本文档的每一张验证表
都有三列，且三列的含义是固定的：

\begin{center}
\small
\begin{tabular}{@{}p{0.16\textwidth} p{0.72\textwidth}@{}}
\toprule
\textbf{列} & \textbf{含义与纪律}\\
\midrule
基准 & 检验的对象。措辞要能回答"若它错了，什么会跟着错"\\
判据 & \textbf{测试断言的门槛}。它是回归门禁，不随版本漂移\\
实测 & \textbf{编写时观测到的值}。它带有效数字、可回查，但会随版本演化；
       引用时必须连同模块版本一并引用\\
\bottomrule
\end{tabular}
\end{center}

\noindent
纪律的另一半是\textbf{什么不能进"实测"列}：仅有趋势正确、
仅有一阶吻合、或无独立基准可对照者，一律进 §x.6 的"适用边界"。
本文档第 \ref{ch:transport_dd} 章 T1 基准中"未做"的一项
（与一维 SP 的带边差）就是这条纪律的应用：它在\emph{原理上}
不可达，因此被写成诚实结论，而不是待办事项。

\cleardoublepage
